\itemhead{4.1 思考 · 未编号题组（教材p.98）}{XJ-B1-C04-TH-P098-Y1358}
\label{XJ-B1-C04-TH-P098-Y1358}
{\footnotesize 湘教版教材p.98 / PDF第105页 / 4.1 思考 / 未编号题组}
\field{原题}{当 $a>0$，$x$ 是任意有理数时，$a^x$ 都有意义。在实际问题中，指数 $x$ 还可以是无理数。那么，正数的无理数指数幂应该怎样定义呢？}
\field{参考答案}{对无理数 $x$，用逐步逼近 $x$ 的有理数指数对应的幂值逼近一个确定正数，将它定义为 $a^x$（$a>0$）。例如按小数逼近指数；这里只作教材层面的连续扩充解释，不要求极限的形式化证明。}
{\small \field{关联知识}{P02 有理数、实数指数幂与运算（教材帮p.144、145 / PDF153、154）}
\field{方法与易错点}{统一为指数幂，分开处理系数与各底数的指数；最终复查定义条件及数值精度。 易错：负指数误作负数；幂的幂指数相加；约分丢禁点；过早四舍五入。}
\field{答案依据}{答案册未提供本栏目的独立对应答案；依据原题、教材概念补写参考回应。开放性总结不设唯一数值答案。}}
